% resumo em português
\begin{resumo}
\noindent
No Instituto Federal Catarinense (IFC), cada componente curricular possui um plano de ensino que deve conter os elementos exigidos pela Organização Didática da instituição, e a conferência desses documentos é feita manualmente pela coordenação e pelo Núcleo Docente Estruturante (NDE), tarefa repetitiva e sujeita a inconsistências entre avaliadores. Este trabalho apresenta o desenvolvimento e a avaliação de um sistema semiautomático de apoio à validação de planos de ensino do curso de Ciência da Computação do IFC -- Campus Videira. A solução, construída sobre a plataforma de automação n8n e executada em contêineres Docker, obtém os planos em PDF de uma pasta do Google Drive, extrai e segmenta o texto em seções, avalia doze critérios com um modelo de linguagem de grande escala (LLM) aberto, o Qwen2.5 de 7 bilhões de parâmetros executado localmente pelo Ollama, e registra o resultado de cada critério (OK, INCOMPLETO ou AUSENTE) em uma planilha consultada por um painel \textit{web}. Para reduzir alucinações, a arquitetura combina o LLM com saída estruturada por esquema JSON e com regras determinísticas, como a contagem de referências bibliográficas por expressões regulares. O desenvolvimento seguiu um processo iterativo e incremental, e a avaliação usou os 29 planos de ensino do semestre 2024/1. Ao longo das iterações, a proporção de planos com falha técnica caiu de 86,2\% para zero, e a versão final processou o lote completo em 8 minutos e 54 segundos, com média de 11,5 segundos de inferência por plano em uma GPU de 6 GB. Na execução final, 11 planos (38\%) foram considerados conformes, 3 (10\%) pendentes e 15 (52\%) reprovados, e a ausência do horário de atendimento ao estudante foi a não conformidade mais frequente. No critério verificável por oráculo estrutural, o sistema obteve \textit{recall} de 100\%, precisão de 93,8\% e Kappa de Cohen de 0,93. A análise de erros identificou limitações, como textos provisórios do sistema acadêmico aceitos como conteúdo válido e observações gerais do modelo contraditórias com o resultado por critério, que orientam os trabalhos futuros. Conclui-se que o sistema é viável como ferramenta de triagem que reduz o esforço de conferência, desde que a decisão final permaneça com o avaliador humano.

 \vspace{0.2cm}
\noindent Palavras-chave: Inteligência artificial. Modelos de linguagem. Plano de ensino. Automação de processos. n8n.

\end{resumo}

% resumo em inglês
\begin{resumo}[Abstract]
\begin{otherlanguage*}{english}
\noindent
At the Instituto Federal Catarinense (IFC), every course unit has a teaching plan that must contain the elements required by the institution's academic regulations, and these documents are checked manually by the program coordination and the Structuring Teaching Nucleus (NDE), a repetitive task prone to inconsistencies between reviewers. This work presents the development and evaluation of a semi-automatic system that supports the validation of teaching plans in the Computer Science program at IFC -- Videira Campus. The solution, built on the n8n automation platform and deployed with Docker containers, retrieves the plans as PDF files from a Google Drive folder, extracts and segments their text into sections, assesses twelve criteria with an open large language model (LLM), the 7-billion-parameter Qwen2.5 running locally on Ollama, and stores the result of each criterion (OK, INCOMPLETE or MISSING) in a spreadsheet displayed by a web dashboard. To reduce hallucinations, the architecture combines the LLM with JSON-schema structured output and with deterministic rules, such as counting bibliographic references with regular expressions. Development followed an iterative and incremental process, and the evaluation used the 29 teaching plans of the 2024/1 semester. Across iterations, the share of plans with technical failures dropped from 86.2\% to zero, and the final version processed the whole batch in 8 minutes and 54 seconds, averaging 11.5 seconds of inference per plan on a 6 GB GPU. In the final run, 11 plans (38\%) were compliant, 3 (10\%) pending and 15 (52\%) rejected, and the missing student office hours was the most frequent non-compliance. On the criterion that could be checked by a structural oracle, the system reached 100\% recall, 93.8\% precision and a Cohen's Kappa of 0.93. The error analysis revealed limitations, such as placeholder texts from the academic system being accepted as valid content and general remarks from the model contradicting the per-criterion results, which guide future work. The system proved viable as a screening tool that reduces review effort, provided that the final decision remains with the human reviewer.

\vspace{0.2cm}
\noindent Keywords: Artificial intelligence. Language models. Teaching plan. Process automation. n8n.
\end{otherlanguage*}
\end{resumo}
